\section{Interaction-state resolution in successor measures}
\label{sec:theory}

We now formalize the interaction-state hypothesis behind our analysis.
The exact successor measure remains well defined in the full physical state, but a learned representation or score may fail to distinguish interaction states with different future occupancies.

\subsection{Full-state successor measures under learned representations}

A full-state successor measure can represent the interaction mode since the interaction state is part of the full physical state.
The error arises when function approximation is introduced to learn state features.
Let $y=\Omega(\tilde{s})$ denote the observation used by the policy or value function, and let $\omega(\tilde{s})$ denote the local interaction state, such as whether an object is freely moving, stably held, or slipping.
Two full states can be nearby in observation space while having different interaction states and therefore different future occupancies.
The relevant question is therefore which interaction-state distinctions are lost in the observation, learned representation, or scalar score.
We illustrate this distinction by decomposing the transition dynamics according to interaction state.

\subsection{Interaction-dependent transition dynamics}

To illustrate how interaction state changes future occupancy, consider a binary projection $\omega(\tilde{s})\in\{0,1\}$.
Here, $\omega=0$ denotes a passive or pre-contact state, while $\omega=1$ denotes a state in which the robot action directly affects the object.
Let $r_t$ denote the robot state and $o_t$ the object state.
For illustration, write the object transition as
\begin{equation}
    o_{t+1}
    \sim
    \begin{cases}
        P_0(\cdot\mid o_t), & \omega_t=0,\\
        P_1(\cdot\mid r_t,o_t,a_t), & \omega_t=1.
    \end{cases}
    \label{eq:mode_switch_target}
\end{equation}
This binary projection records whether the robot action directly affects the object; finer interaction states may distinguish stable control, slipping, and loss of control.
When $\omega_t=0$, the object transition does not depend directly on the robot action.
When $\omega_t=1$, it depends on the robot state and action.
Thus, similar robot--object configurations can have different successor measures when their interaction states induce different object dynamics.

\subsection{Control acquisition and task success}
\label{sec:tail_successor_gaps}

In interaction-rich tasks, 
acquiring control of the object separates reaching it from completing the task.
A policy may satisfy a coarse control-acquisition criterion while reaching a state from which retaining and transporting the object remains unlikely.

Let $C\subseteq\widetilde{\mathcal S}$ denote the control-acquisition set and define its first hitting time as $\tau_C=\inf\{t:\tilde{s}_t\in C\}$.
In cube manipulation, states in $C$ may correspond to either a stable grasp or an unstable grasp that is likely to slip.
The event $\tau_C\leq H$ is therefore not sufficient for task success: success also depends on the state from which the policy begins transporting the cube.

For a remaining horizon $k$, define the task-success value from state $\tilde{s}$:
\begin{equation}
    v_k^\pi(\tilde{s};g) = \Pr_\pi \left( \tilde{s}_k\in R\cap G_g \mid \tilde{s}_0=\tilde{s} \right).
    \label{eq:tail_success}
\end{equation}
Conditioning on the first control-acquisition time gives
\begin{equation}
    u_H^\pi(\tilde{s},a;g) = \mathbb{E}_\pi \left[ \mathbf{1}\{\tau_C\leq H\} v_{H-\tau_C}^\pi(\tilde{s}_{\tau_C};g) \mid \tilde{s}_0=\tilde{s},a_0=a \right].
    \label{eq:bridge_decomposition}
\end{equation}
This decomposition separates whether the policy acquires control from whether it completes the task after control has been acquired.

Our experiments show that most performance differences arise after control acquisition.
Policies generally reach and lift the cube, but many fail to retain it or transport it to the goal.
Stable and unstable grasp states can both satisfy the same lifting criterion while assigning future occupancy to different outcomes.
A stable grasp places more future occupancy on retained-cube and goal states, whereas an unstable grasp places more occupancy on dropped-cube or repeated-pickup states.

\subsection{Local successor gaps in learned scores}

For an observation scale $\epsilon>0$ and remaining horizon $k$, define the local successor gap
\begin{equation}
    \Gamma^\pi_{\epsilon,k}(g) = \sup_{\substack{ \tilde{s}_1,\tilde{s}_2\in C\\
    \|\Omega(\tilde{s}_1)-\Omega(\tilde{s}_2)\|\leq\epsilon }}
    \left| v_k^\pi(\tilde{s}_1;g) - v_k^\pi(\tilde{s}_2;g) \right|.
    \label{eq:tail_gap}
\end{equation}
A large $\Gamma^\pi_{\epsilon,k}(g)$ means that nearby states reached after lifting have substantially different task-success values.
Such a gap arises when two states satisfy the same lifting criterion, but only one is likely to retain the cube and reach the goal.

Suppose a stable grasp state $\tilde{s}_s$ and an unstable grasp state $\tilde{s}_u$ produce the same observation,
\[
    \Omega(\tilde{s}_s)=\Omega(\tilde{s}_u)=y,
\]
but have different task-success values:
\[
    v_s=v_k^\pi(\tilde{s}_s;g), \qquad v_u=v_k^\pi(\tilde{s}_u;g).
\]
If $p$ is the probability that observation $y$ corresponds to the stable grasp state, then the optimal observation-only predictor under squared error is
\begin{equation}
    \bar v(y)=p v_s+(1-p)v_u.
\end{equation}
Its errors on the two modes are
\begin{equation}
    |v_s-\bar v|=(1-p)|v_s-v_u|, \qquad |v_u-\bar v|=p|v_s-v_u|.
\end{equation}
Thus, if stable and unstable grasp states receive the same observation or representation, an observation-based value estimate must underestimate the stable state, overestimate the unstable state, or both.
Overestimating unstable grasp states can cause the actor to prefer actions that lift the cube but do not maintain control during transport.

This information loss may occur either in the learned representation or in the scalar readout.
A bilinear score can express the distinction only if the representation encodes the relevant interaction-state information and the task-conditioned readout uses it.
Our empirical probes therefore ask whether the learned representation contains information about task success from lifted-cube states and whether the scalar score preserves that information.

\subsection{Instantiation in cube manipulation}

We study this hypothesis in the state-based \ogbench{} \task{} environment.
Here, the relevant interaction state describes whether the cube is freely moving, stably held, slipping, or no longer controlled.
The observation contains the gripper pose, cube pose, gripper opening, and contact indicators, but the distinction between stable and unstable grasps is distributed across these state variables.
Consequently, nearby robot--cube configurations can differ in whether the policy will retain the cube during transport.

For goal $g$, we use the task-specific event
\begin{equation}
    \varphi_g(\tilde{s}_{0:H}) = \mathbf{1}\{
        \tau_A\leq H,\;
        \tau_C\leq H,\;
        \tilde{s}_H\in R,\;
        \tilde{s}_H\in G_g
    \},
\end{equation}
where $A$ is the reaching set, $C$ is the control-acquisition set, $R$ is the final-retention set, and $G_g$ is the goal set.
The corresponding task-success probability is
\begin{equation}
    u^\pi_H(\tilde{s},a;g) = \Pr_\pi
    \left(
        \tau_A\leq H,\;
        \tau_C\leq H,\;
        \tilde{s}_H\in M,\;
        \tilde{s}_H\in G_g
        \mid
        \tilde{s}_0=\tilde{s},a_0=a
    \right).
\end{equation}

The four criteria distinguish reaching the cube, lifting it, retaining it at the end of the episode, and transporting it to the goal.
A final-retention failure means that the policy satisfied the lifting criterion but did not end the episode with the cube above the retention threshold.